\documentclass{article} %210 mm × 297 mm
\usepackage[utf8]{inputenc}
\usepackage[a4paper,left=1.5cm, right=1.5cm, top=2cm, bottom=2cm]{geometry}
\usepackage{graphicx}
%\usepackage{blindtext}
\usepackage{multirow}
\usepackage{mhchem}
\usepackage{url}
\usepackage{subfigure}
\usepackage{amsfonts,latexsym} % para tener disponibilidad de diversos simbolos
\usepackage{enumerate}
%\usepackage{booktabs}
\usepackage{float}
%\usepackage{threeparttable}
\usepackage{array,colortbl}
%\usepackage{ifpdf}
%\usepackage{rotating}
\usepackage{cite}
\usepackage{stfloats}
\usepackage{url}
\usepackage{listings}
\usepackage{fancyhdr}

\usepackage{color}
%\usepackage{hyperref}
%\usepackage{wrapfig}
\usepackage{array}
%\usepackage{multirow}
%\usepackage{adjustbox}
%\usepackage{nccmath}
\usepackage[dvipsnames]{xcolor}


\newcommand{\titulo}{Abgabe 3; MfN II.}
\newcommand{\nombre}{Liliana Sanfilippo}
\newcommand{\FCE}{\textbf{SoSe 2025}}

 

\usepackage{fancyhdr}
\pagestyle{fancy}
\fancyhead{}
\fancyfoot{}
\fancyhead[LO]{\small\nombre}
\fancyhead[RE]{\small\titulo}
\fancyhead[LO]{\small\FCE}
\fancyfoot[RO,LE] {\thepage}

\setcounter{page}{1} %Número de la página, la editorial lo cambiará

\newcommand{\blue}[1]{\textcolor{blue}{#1}}
\newcommand{\ora}[1]{\textcolor{orange}{#1}}


\usepackage[onehalfspacing]{setspace}
\begin{document}


\begin{tabular}{p{12cm}}
{{\Large\bf\titulo}}\\
\vspace{0.2cm}\\
 {\large\nombre}\\
\end{tabular}
\vspace{0.2cm}\\
\section{Aufgabe 3.1}
\subsection*{Matrix C}
\[
C = \begin{pmatrix}
3&1&1\\
2&4&2\\
1&1&3
\end{pmatrix}
\]
\[
C -\lambda I = \begin{pmatrix}
3-\lambda&1&1\\
2&4-\lambda&2\\
1&1&3-\lambda
\end{pmatrix}
\]
\[
det (C) = (-\lambda+3)\cdot(-\lambda+4)\cdot(-\lambda+3)+1\cdot2\cdot1+1\cdot2\cdot1-1\cdot(-\lambda+4)\cdot1-1\cdot2\cdot(-\lambda+3)-(-\lambda+3)\cdot2\cdot1 = -\lambda^3+10\cdot\lambda^2-28\cdot\lambda+24
\]
\[
\Rightarrow \chi_C = -\lambda^3+10\cdot\lambda^2-28\cdot\lambda+24
\]
\begin{align*}
    -\lambda^3+10\cdot\lambda^2-28\cdot\lambda+24 & = -(\lambda-2) \cdot (\lambda^2 -8\lambda+12) \\
    & = -(\lambda-2) \cdot (\lambda-2) \cdot (\lambda-6)
\end{align*}
\[
\lambda_1 = 2, \quad \lambda_2 = 2, \quad \lambda_3 = 6
\]
\subsubsection*{Für Eigenwert 2:}
\[
C - 2 I = \begin{pmatrix}
3-2&1&1\\
2&4-2&2\\
1&1&3-2
\end{pmatrix}
 = \begin{pmatrix}
1&1&1\\
2&2&2\\
1&1&1
\end{pmatrix}
\]
Setze $(A - 2)v = 0$.
\[
\begin{pmatrix}
1&1&1\\
2&2&2\\
1&1&1
\end{pmatrix} \cdot \begin{pmatrix}
a \\ b\\c
\end{pmatrix} = \begin{pmatrix}
0\\0\\0
\end{pmatrix}
\]
Gleichungssystem lösen (für Zeile 1, da alle davon linear abhängig). 
\begin{align*}
    0 & = a+b+c \\
    \Rightarrow a & = -b-c 
\end{align*}
EV hat Form 
\[
\begin{pmatrix}
-b-c\\b\\c
\end{pmatrix} = b \cdot \begin{pmatrix}
-1\\1\\0
\end{pmatrix} + c \cdot \begin{pmatrix}
-1\\0\\1
\end{pmatrix}
\]


\subsubsection*{Für Eigenwert 6:}
\[
C - 2 I = \begin{pmatrix}
3-6&1&1\\
2&4-6&2\\
1&1&3-6
\end{pmatrix}
 = \begin{pmatrix}
-3&1&1\\
2&-2&2\\
1&1&-3
\end{pmatrix}
\]
Setze $(A - 6)v = 0$.
\[
\begin{pmatrix}
-3&1&1\\
2&-2&2\\
1&1&-3
\end{pmatrix} \cdot \begin{pmatrix}
a \\ b\\c
\end{pmatrix} = \begin{pmatrix}
0\\0\\0
\end{pmatrix}
\]
Gleichungssystem lösen 
\[
\begin{array}{cccc}
 -3a&+b&+c&=0\\
2a&-2b&+2c&=0\\
a&+b&-3c&=0
\end{array}
\]

\begin{align*}
    0 & = 2a-2b+2c\\
    \Rightarrow   0 & = a-b+c\\
    \Rightarrow   b & = a+c\\
\end{align*}
\[
 -3a+b+c  = -3a+a+c+c = -2a+2c = 0
\]
\[
\Rightarrow c = a
\]
EV hat Form:
\[
 \begin{pmatrix}
a\\a+c\\a
\end{pmatrix}
\]
Also sind die Eigenvektoren:
\[
v_1 = \begin{pmatrix}
-1\\1\\0
\end{pmatrix}, \quad v_2 = \begin{pmatrix}
-1\\0\\1
\end{pmatrix}, \quad v_3 =  \begin{pmatrix}
1\\2\\1
\end{pmatrix}
\]

\subsection*{Matrix D}
\[
D = \begin{pmatrix}
2-i&0&0\\
i&-3&1-i\\
1&0&2+i
\end{pmatrix}
\]
\[
D - \lambda I = \begin{pmatrix}
2-i- \lambda&0&0\\
i&-3- \lambda&1-i\\
1&0&2+i- \lambda
\end{pmatrix}
\]
\begin{align*}
    det(D) & = (-\lambda+(2-i))\cdot(-\lambda-3)\cdot(-\lambda+(2+i)) \\
    & +0\cdot(1-i)\cdot1+0\cdot i\cdot0-1\cdot(-\lambda-3)\cdot0-0\cdot(1-i)\cdot(-\lambda+(2-i)) \\
    & -(-\lambda+(2+i))\cdot i\cdot0 \\
   & = (-\lambda+(2-i))\cdot(-\lambda-3)\cdot(-\lambda+(2+i)) +0+0-0-0 -0 \\
   & = (-\lambda+(2-i))\cdot(-\lambda-3)\cdot(-\lambda+(2+i))\\
      & = (-\lambda+(2-i))\cdot(-\lambda-3)\cdot(-\lambda+(2+i))\\
      & = ((2-\lambda)^2 - i^2) \cdot (-\lambda-3) \\
      & = ((2-\lambda)^2 +1) \cdot (-\lambda-3) \\
    & = (5-4\lambda+ \lambda^2) \cdot (-\lambda-3) \\
    & = -\lambda^3 + \lambda^2 + 7\lambda -15 
\end{align*}
\[
\chi_D = -\lambda^3 + \lambda^2 + 7\lambda -15 
\]
bzw. 
\[
\chi_D = \lambda^3 - \lambda^2 - 7\lambda +15 
\]
Nullstellen
\[
\lambda^3 - \lambda^2 - 7\lambda +15  = (\lambda +3) \cdot (\lambda^2 - 4 \lambda +5)
\]
und Mitternachtsformel fpr $\lambda^2 - 4 \lambda +5$
\[
\lambda_{2,3} = \frac{4 \pm \sqrt{(-4)^2 - 4 \cdot 1 \cdot 5 }}{2 \cdot 1} = \frac{4 \pm \sqrt{16 -20} }{2} = \frac{4 \pm \sqrt{-4}}{2} = \frac{4 \pm 2i}{2} = 2 \pm i
\]
Also sind die Nullstellen: 
\[
\Rightarrow \lambda_1 = -3, \quad \lambda_2 = 2+i, \quad \lambda_3 = 2-i
\]
\subsubsection*{Für Eigenwert -3:}
\[
C + 3 I = \begin{pmatrix}
3+3&1&1\\
2&4+3&2\\
1&1&3+3
\end{pmatrix}
 = \begin{pmatrix}
6&1&1\\
2&7&2\\
1&1&6
\end{pmatrix}
\]
Setze $(A + 3)v = 0$.
\[
\begin{pmatrix}
6&1&1\\
2&7&2\\
1&1&6
\end{pmatrix} \cdot \begin{pmatrix}
a \\ b\\c
\end{pmatrix} = \begin{pmatrix}
0\\0\\0
\end{pmatrix}
\]
Lineares Gleichungssystem lösen: 
\begin{align*}
            0 & =  6a+b+c \\
\Rightarrow c & = -6a-b \\
            0 & = a+b+6c \\
              & = a+b+6\cdot(-6a-b) \\
              & = a+b-36a-6b\\
              & = -35a-5b \\
\Rightarrow b & = -7a \\ 
\Rightarrow c & = -6a+7a = a 
\end{align*}
EV hat die Form
\[
\begin{pmatrix}
1 \\ -7\\1
\end{pmatrix}
\]
\subsubsection*{Für Eigenwert 2+i:}
\[
C  (-2+i) I = \begin{pmatrix}
3-(2+i)&1&1\\
2&4-(2+i)&2\\
1&1&3-(2+i)
\end{pmatrix}
 = \begin{pmatrix}
1-i&1&1\\
2&2-i&2\\
1&1&1-i
\end{pmatrix}
\]
Setze $(A -(2+i) )v = 0$.
\[
\begin{pmatrix}
1-i&1&1\\
2&2-i&2\\
1&1&1-i
\end{pmatrix} \cdot \begin{pmatrix}
a \\ b\\c
\end{pmatrix} = \begin{pmatrix}
0\\0\\0
\end{pmatrix}
\]
Lineares Gleichungssystem lösen:
\begin{align*}
            0 & = a(1-i)+b+c \\
\Rightarrow a & = \frac{-b-c}{1-i} \\
            0 & = a+b+(1-i)c \\
\Rightarrow a & = -b-(1-i)c\\
              & = -b-c+ic  \\ \\
\frac{-b-c}{1-i} & =  -b-c+ic  \\
\Rightarrow -b-c & = (-b-c+ic) (1-i) \\
\Rightarrow-c & = bi+2ic \\
              & = i(b+2c) \\
\Rightarrow \frac{-c}{i} & = b+2c \\
\Rightarrow b & = \frac{-c}{i} -2c  \\
              & = \frac{-c}{i} \cdot \frac{i}{i} -2c \\
              & = \frac{-ci}{i^2} -2c \\
              & = \frac{-ci}{-1} -2c \\
              & = ci -2c \\ \\
\Rightarrow a & = -(ci -2c)-(1-i)c\\
              & = -ci + 2c +ci -c \\
              & = c \\ \\
\end{align*}
EV hat die Form
\[
\begin{pmatrix}
1 \\i-2 \\1
\end{pmatrix}
\]
\subsubsection*{Für Eigenwert 2-i:}
\[
C  (-2-i) I = \begin{pmatrix}
3-(2-i)&1&1\\
2&4-(2-i)&2\\
1&1&3-(2-i)
\end{pmatrix}
 = \begin{pmatrix}
1+i&1&1\\
2&2+i&2\\
1&1&1+i
\end{pmatrix}
\]
Setze $(A -(2-i) )v = 0$.
\[
\begin{pmatrix}
1+i&1&1\\
2&2+i&2\\
1&1&1+i
\end{pmatrix} \cdot \begin{pmatrix}
a \\ b\\c
\end{pmatrix} = \begin{pmatrix}
0\\0\\0
\end{pmatrix}
\]
Lineares Gleichungssystem lösen:
\begin{align*}
            0 & = a(1+i) +b+c \\
\Rightarrow b & = -c -a(1+i) \\ \\ 
%            0 & = 2a(1+i) +b(2+i) + 2c \\
%              & = 2a(1+i) + (-c -a(1+i))(2+i) + 2c \\
%              & = 2a + 2ai + (-c -a-ai)(2+i) + 2c \\
%              & = 2a + 2ai - c(2+i) -a(2+i) -ai(2+i) +2c \\
%              & = 2a + 2ai - 2c -ci -2a -ai -2ai -ai^2 +2c \\
%              & =  -ci -ai - a\cdot(-1) \\
%              & =  -ci -ai +a\\
%              & = a + i(-c-a) \\ \\
            0 & = a + b + (1+i)c \\
              & = a + -c -a -ai  + c+ ci \\
              & = -ai + ci \\
\Rightarrow a & = c \\
\Rightarrow b & = -a -a(1+i) \\
              & = -a -a -ai \\
              & = -2a -ai
\end{align*}
EV hat Form: 
\[
\begin{pmatrix}
1 \\ -2-i\\1
\end{pmatrix}
\]
Also sind die Eigenvektoren:
\[
v_1 = \begin{pmatrix}
1 \\ -7\\1
\end{pmatrix}, \quad v_2 = \begin{pmatrix}
1 \\i-2 \\1
\end{pmatrix}, \quad v_3 =  \begin{pmatrix}
1 \\ -2-i\\1
\end{pmatrix}
\]
\subsection*{Matrix E}
\[
 E = \begin{pmatrix}
 1 &-2&1 \\
 1&1&3\\
 0&-1&-2
 \end{pmatrix}
\]
\[
 E -\lambda I= \begin{pmatrix}
 1-\lambda&-2&1 \\
 1&1-\lambda&3&\\
 0&-1&-2-\lambda
 \end{pmatrix}
\]
\begin{align*}
    det(E) & =  (-\lambda+1)\cdot(-\lambda+1)\cdot(-\lambda-2)+
    \ora{(-2)\cdot3\cdot0}
    + \blue{1\cdot1\cdot(-1)}
    -\ora{0\cdot(-\lambda+1)\cdot1}
    -\blue{(-1)}\cdot3\cdot(-\lambda+1)
    -(-\lambda-2)\cdot\blue{1}\cdot(-2)  \\
& =  (-\lambda+1)\cdot(-\lambda+1)\cdot(-\lambda-2)
    + \blue{(-1)}
    -(-3)\cdot(-\lambda+1)
    -(-\lambda-2)\cdot(-2)  \\
    & = (-\lambda+1)\cdot(-\lambda+1)\cdot(-\lambda-2)
    -1
    -(-3)\cdot(-\lambda+1)
    -(-\lambda-2)\cdot(-2)  \\
     & = -\lambda^3 + 3 \lambda -2
    -1
    -3\lambda+3
    -2\lambda -4\\
    & = -\lambda^3 -2\lambda -4\\
\end{align*}
\[
\chi_E = -\lambda^3 -2\lambda -4
\]
bzw. 
\[
\chi_E = \lambda^3 +2\lambda +4
\]
Nullstellen finden mit Cardano
\[ 0  = \lambda^3 +2\lambda +4\]
Wolfram Alpha löst auf (keine Ahnung wie das händisch geht, frage ich im nächsten Tutorium) als:
\begin{align*}
 \lambda_1 &  \approx -1,1795 \\
  \lambda_2 &  \approx 0,5898 - 1,7445 i \\
   \lambda_3 &  \approx 0,5898 - 1,7445 i
\end{align*}
Wenn man $E$ als reelle $3\times3$-Matrix auffasst, findet man nur einen Eigenwert, die anderen beiden sind weiterhin komplex. Ohne Eigenwerte keine Eigenvektoren. 
\section{Aufgabe 3.2}
Wenn alle Einträge in $A$ reell sind, können keine $i$s drin sein. Daher ist $A = \bar A$. \\
Daher gilt auch
\[ \bar A\cdot v  = A\cdot v = \lambda \cdot v\]
bzw.
\[\overline{A\cdot v} = \overline{\lambda \cdot v} \Rightarrow \bar A\cdot \bar v = \bar\lambda \cdot \bar v \Rightarrow  A \cdot \bar v = \bar\lambda \cdot \bar v\]
Und damit ist $\bar v$ ein Eigenvektor zum Eigenwert $\bar \lambda$ für $A$. 
\section{Aufgabe 3.3}
\[
A = \begin{pmatrix}
0&\cdots&0&-\alpha_0\\
1&&0&-\alpha_1\\
& \ddots & & \vdots\\
0&&1&-\alpha_{n-1}
\end{pmatrix}
\]
Wikipedia sagt $A$ ist die Begleitmatrix eines normierten Polynoms $n$-ten Grades $t^n + \alpha_n t ^{n-1} + \cdots + \alpha_1 t + \alpha_0 $.  \\
Wir wissen aus der Vorlesung das charackteristische Polynom von $A$ ist $\chi_A(t) = det(A - t I)$. \\ \\
z.z. $det(A - t I) = (-1)^n \cdot (t^n + \alpha_{n-1} t ^{n-1} + \cdots + \alpha_1 t + \alpha_0)$
\[
A - t I = \begin{pmatrix}
-t&\cdots&0&-\alpha_0\\
1&-t&0&-\alpha_1\\
& \ddots &\ddots & \vdots\\
0&&1&-\alpha_{n-1} -t
\end{pmatrix}
\]
Leibnitz Entwickung entlang der letzten Zeile. 
\begin{align*}
    det(A - t I) &= \sum^n_{i = 1} a_{in} \cdot (-1)^{i+n} \cdot det(A'_{in}) \\
    &=  \sum^n_{i = 1} a_{in} \cdot (-1)^{n} \cdot (-1)^{i} \cdot det(A'_{in}) 
  %  &= (-1)^{n} \cdot \sum^n_{i = 1} a_{in} \cdot  (-1)^{i} \cdot det(A'_{in})\\
\end{align*}
Diese Summe hat die Form
\[
-\alpha_{0} \cdot (-1)^{1} \cdot det(A'_{1n})  
+ (-\alpha_{1}) \cdot (-1)^{2} \cdot det(A'_{2n}) 
+ ... 
+ (-\alpha_{n-2}) \cdot (-1)^{n-1} \cdot det(A'_{(n-1)n}) 
+ (-\alpha_{n-1}-t) \cdot (-1)^{n} \cdot det(A'_{nn}) 
\]
Dabei hat $A'_{nn}$ die gleiche Form wie $A$, nur ist eine $(n-1)\times(n-1)$-Matrix. Die restlichen $A'_{in}$ haben die Form 
\[
A'_{in} = \begin{pmatrix}
-t & 0 & 0 & 0 & 0 & 0 & 0 & 0 \\
1 & -t & 0 & 0 & 0 & 0 & 0 & 0  \\
0 &\ddots & \ddots & \ddots &\ddots &\ddots &\ddots &0  \\  
0 & 0 & 1 & -t & 0 & 0 & 0 & 0  \\
0 & 0 & 0 & 0 & 1 & -t & 0 & 0  \\
0 &\ddots & \ddots & \ddots &\ddots &\ddots &\ddots  &0\\
0 & 0 & 0 & 0 & 0 & 0 & 1 & -t \\
\end{pmatrix}
\]
Dabei ist das obere linke Viertel eine Matrix der Form 
\[
O=\begin{pmatrix}
-t & 0 & 0 &0 \\
1 & -t & 0 & 0 \\
0 &\ddots & \ddots & 0  \\  
0 & 0 & 1 & -t 
\end{pmatrix}
\]
Und das untere rechte Viertel eine Matrix der Form 
\[
U= \begin{pmatrix}
 1 & -t & 0 & 0  \\
 0 & 1 & -t & 0  \\
0&\ddots &\ddots  &-t\\
0 & 0 &0 & 1   \\
\end{pmatrix}
\]
Also ist jedes $A'_{in}$ eine Blockmatrix und die Determinanten davon kann man berechnen als Produkt der Determinanten der beiden Matrix-Blöcke. \\
Da $U, O$ $(n-1)\times(n-1)$Dreiecksmatrizen sind, gilt hier $det(U) = 1^{n-1} = 1$ und $det(O) = (-t)^{n-1}$.
Also
\[
det(A'_{in}) = det(U) \cdot det(O) = 1 \cdot (-t)^{n-1} = (-t)^{n-1} \quad für\ i \in \{1,...,n-1\}
\] 
und für $A'_{nn}$ gilt ebenfalls $det(A'_{nn}) = (-t)^{n-1}$. \\
Daher 
\begin{align*}
    det(A - t I) & =  \sum^n_{i = 1} a_{in} \cdot (-1)^{n} \cdot (-1)^{i} \cdot det(A'_{in})  \\
    & =  \sum^n_{i = 1} a_{in} \cdot (-1)^{n} \cdot (-1)^{i} \cdot (-t)^{n-1}  \\
    & = (-1)^{2n} \cdot (-\alpha_{n-1} -t)  \cdot (-t)^{n-1} +  \sum^{n-1}_{i = 1} a_{in} \cdot (-1)^{n} \cdot (-1)^{i} \cdot (-t)^{n-1}  \\
     & = (-\alpha_{n-1} -t)  \cdot (-t)^{n-1} +  \sum^{n-1}_{i = 1} a_{in} \cdot (-1)^{n} \cdot (-1)^{i} \cdot (-t)^{n-1}  \\
       & = t^3 + \alpha_{n-1}\cdot t^{n-1} +  \sum^{n-1}_{i = 1} a_{in} \cdot (-1)^{n} \cdot (-1)^{i} \cdot (-t)^{n-1}  \\
        & = t^3 + \alpha_{n-1}\cdot t^{n-1} + \alpha_{n-2}t^{n-2} + \sum^{n-2}_{i = 1} a_{in} \cdot (-1)^{n} \cdot (-1)^{i} \cdot (-t)^{n-1}  \\
       & \quad  \vdots
\end{align*}
\end{document}